\chapter{Fazit und Ausblick}
Dieses Kapitel fasst die Ergebnisse der Arbeit zusammen, bewertet den erreichten Stand und skizziert mögliche Weiterentwicklungen sowie Einsatzmöglichkeiten des Prototyps.

\section{Zusammenfassung der Ergebnisse}
In dieser Arbeit wurde ein universeller Ausleseprototyp für digitale Tachographen konzipiert, aufgebaut und validiert. Die Hardwarearchitektur mit ESP32-S3, Schutzstufen und Pegelwandlern ermöglicht einen robusten Betrieb im 24-V-Bordnetz. Die Firmware implementiert die relevanten Kommunikationsprotokolle und stellt den Download von Massenspeicher- und Fahrerkartendaten bereit. Ergänzend wurde eine plattformübergreifende Anwendung entwickelt, die die ausgelesenen Daten strukturiert visualisiert und grundlegende Auswertungen ermöglicht. Insgesamt zeigt der Prototyp die technische Machbarkeit und bildet eine belastbare Basis für weitere Verbesserungen.

\section{Potenzielle Weiterentwicklungen}
Für die Weiterentwicklung stehen insbesondere die Stabilisierung der Tachographenkommunikation, erweiterte Gerätekompatibilität und zusätzliche Tests im Fahrzeugumfeld im Fokus. Darüber hinaus bieten sich eine weitergehende Automatisierung der Ausleseprozesse, eine verbesserte Datenvalidierung sowie eine sichere Cloud-Anbindung an. Auf Softwareseite sind eine erweiterte Nutzerverwaltung, Suchfunktionen und detailliertere Auswertungen denkbare Schritte, um den praktischen Nutzen der Anwendung zu erhöhen.
Geplant sind ein Nutzerkonto mit Anbindung an eine Datenbank zur Zuordnung und zum Abruf von \code{.ddd}-Dateien sowie eine Suchfunktion für gezielte Abfragen. Zusätzlich soll das manuelle Hinzufügen von Tachographendateien ohne ESP32 möglich werden. Weitere Punkte sind Designanpassungen, eine integrierte Bedienungsanleitung und zusätzliche Hilfetexte.
